\chapter{CONCLUSION AND RECOMMENDATION}

\section{Conclusions}
The comparison of the four analytical support conditions showed that SSI had a measurable influence on the seismic response of the selected hospital prototype. The first-mode period increased from 0.984~s for fixed support to 1.147~s for the soft-soil model, while the 90 percent cumulative modal-mass criterion was achieved by Mode~6 in both principal directions for every model.

Response-spectrum displacement increased as support flexibility increased. RSx ULS roof displacement increased from 87.224~mm to 162.756~mm and RSx SLS roof displacement from 83.716~mm to 156.210~mm between fixed and soft support. RSy ULS roof displacement increased from 71.758~mm to 120.531~mm, while RSy SLS roof displacement increased from 88.944~mm to 170.330~mm.

Inter-storey drift followed the same support-stiffness trend. Maximum RSx ULS drift increased from 0.009385 to 0.013882 and maximum RSx SLS drift from 0.009007 to 0.013324. The corresponding maximum RSy ULS and SLS drifts increased from 0.004960 to 0.007677 and from 0.004761 to 0.007369, respectively.

Force response did not vary uniformly with soil flexibility. RSx ULS base shear increased from 21,518.409~kN to 24,528.262~kN and RSx SLS base shear from 20,652.947~kN to 23,541.745~kN between fixed and soft support. Lower-storey RSx shear also increased toward the soft-soil case. Beam V2 and M3 demands increased substantially, whereas maximum column axial force changed only slightly and column M3 increased. These results indicate redistribution of demand rather than a uniformly beneficial or adverse SSI effect.

Foundation response followed the same support-stiffness ordering. The ETABS settlement series increased from 0.000~mm for the fixed-base reference to 9.209~mm for the soft-soil model, consistent with the decreasing translational and rotational spring stiffness assigned to the flexible-base cases.

The fixed-base-normalized linear sensitivity comparison gave a consistent soil-condition ranking of hard $<$ medium $<$ soft for fundamental period, RSx roof displacement, RSx inter-storey drift, RSx base shear, beam moment, and column moment. The soft-soil case produced the largest relative changes in all selected quantities, including approximately 95.3 percent in maximum beam $|M3|$, 86.6 percent in RSx ULS roof displacement, and 47.9 percent in RSx ULS maximum drift. Settlement likewise reached its largest value, 9.209~mm, for the soft-soil model.

The SAP2000 pushover analyses provided the nonlinear static capacity comparison. The capacity curves are shown only over the displacement range reached in each analysis. Eight X- and Y-direction performance points for the fixed, hard-soil, medium-soil, and soft-soil models were compared within the FEMA~356 target-displacement framework. The X-direction results showed lower base shear with increasing support flexibility while the performance displacement remained approximately 0.12~m; the Y-direction results showed increasing displacement from fixed to flexible support.

The fragility results from the twelve ground motions indicated progressively lower mean damage-state threshold PGA as support flexibility increased. The IO mean threshold decreased from 0.2828~$g$ for fixed support to 0.1747~$g$ for soft support, with the same fixed--hard--medium--soft ordering for OP, DC, LS, and CP. These findings remain preliminary because the available nonlinear results do not specify the response direction and the Barpark/Barpak-labelled record could not be independently verified.

\section{Recommendations}
SSI should be considered in the seismic assessment of hospital buildings founded on flexible soils, particularly where deformation-sensitive performance is important. Displacement, inter-storey drift, foundation response, and member-force redistribution should be evaluated together rather than relying on base shear or a fixed-base idealization alone.

Site-specific geotechnical information should be used where available to define soil stiffness and foundation response. Sensitivity to foundation geometry and alternative SSI representations should also be considered when the expected foundation flexibility is significant.

For nonlinear static assessment, the FEMA~356 target-displacement inputs and calculation record should be retained with the structural-analysis results so that reported performance coordinates can be independently reproduced. Fragility studies should likewise retain the response direction and ground-motion identification needed to support directional interpretation and repeatable analysis.

\section{Limitations of the Study}
The conclusions apply to the selected hospital prototype, the adopted foundation-spring idealization, and the analysed support conditions. The spring representation does not substitute for a site-specific geotechnical investigation. The available nonlinear time-history results do not specify the response direction, and the Barpark/Barpak-labelled record could not be independently verified; consequently, the fragility results are preliminary. Threshold PGAs above the directly analysed 0.1--0.5~$g$ range are extrapolated. The FEMA~356 performance points were not independently recalculated because the intermediate target-displacement coefficient calculation sheet was unavailable. For the 2015 Gorkha earthquake record, only the selected acceleration component with PGA 0.0876~$g$ and a 0.01~s sampling interval was used for the PGA comparison and acceleration time history.
